Spatial Branch and Bound (BB) is a method to find global optima of nonlinear mathematical programs by means of splitting the feasible region in a tree-like pattern and constructing locally-valid cuts from its constraints. BB maintains upper and lower bounds on the global optimum, and the gap between them serves as a proof on how suboptimal the current solution could be. When the gap is smaller than a user-specified threshold, the algorithm stops.

\begin{figure}
    \centering
    \includegraphics[]{babtree.pdf}
    \caption{An example of a Spatial Branch and Bound tree of a nonconvex program minimizing the variable $y$ over $x \in [0, 4]$ such that it lies above the indicated quadratic constraints. The infeasible region is shaded in red, and the locally-valid relaxations via secant cuts are hatched in black.}
    \label{fig:babtree}
\end{figure}

Each step of the algorithm is a branch associated with a set of local bounds on the variables. At each step, the method constructs a convex relaxation for each constraint separately; the relaxed constraints are then solved, e.g. using linear programming. Solutions of the relaxations are lower-bounds on the optimum within the branch and any of its child branches. The worst-case lower-bound in any currently active branch is then a lower bound on the global optimum of the original problem. To improve the lower-bounds, a variable is selected to branch on and its bounds are split into several sub-regions. At the same time, any admissible point for the original problem is an upper-bound on the global optimum. While the program could branch deep enough for the local relaxation to yield an admissible point, it is usual to use heuristics to accelerate the search for feasible points. Any branch where the relaxed subproblem gives a value worse than an admissible solution can be pruned away. \cref{fig:babtree} illustrates how BB uses progressively tighter locally-valid cuts to improve the approximation of a concave feasible region defined by quadratic inequalities.

In the rest of the paper, we will show that BB can be performant and scale to large problem sizes, on top of being widely applicable. Even in cases where established techniques exist, such as Lasserre's hierarchies applied to polynomial optimization, BB methods can be competitive. In fact, recent results \cite{dey2021branch} indicate that considering random instances of some problems results in polynomial time complexity for BB methods.

\subsection{Convex Relaxations}
During the bounding phase of BB, the algorithm constructs a relaxation of the original problem, transforming it into a form that is practically solvable. This likely means relaxing the original problem into a convex one. This process relies on the ability to create reasonable relaxations of the constraints. The literature provides tight convex relaxations for common nonconvex terms, including bilinear and fractional terms, exponentiation, and wholly convex or concave functions. Although tight relaxations may not be generically available, the existing set of relaxations is sufficient to address all algebraic functions \cite{SMITH1999}.

Algebraic functions can be decomposed into a hierarchy of binary operations, each corresponding to one of the aforementioned functions. Consequently, all algebraic functions can be solved using BB by introducing auxiliary variables and appropriate constraints. An implementation of such an algorithm is currently available in the \href{https://github.com/coin-or/Couenne}{Couenne} solver \cite{CoinOr}. As an example, to address the multilinear constraint $xyz=0$, we could introduce a lifting variable $\alpha=xy$ and rewrite the original constraint as $\alpha z=0$. The result are two constraints for which tight convex relaxations exist.

Symbolic reformulations increase the dimensionality of the problem, which is a significant concern given that the worst-case complexity of the branch-and-bound algorithm is exponential in the number of variables. However, in \cref{ch:ik}, we will demonstrate that these transformations can be competitive when compared to Lasserre's hierarchies, particularly for inverse kinematics of generic redundant manipulators. The polynomial problems arising from inverse kinematics constraints  involve more than fourteen variables and degrees greater than four.

\subsection{Global Optimality Certificates}
In comparison to the obvious certificate of global optimality obtained when a feasible solution is found by an outer approximation method such as the Lasserre's hierarchies described in \cref{sec:momsos}, proving the optimality for the spatial Branch and Bound is less direct.
While solvers could return a certificate proving that the solution found is globally optimal within a predefined margin of error, this would reveal the cuts taken by the solver. In the rest of this manuscript, we will focus on the Gurobi solver, which uses a closed-source variant of the spatial BB method and does not return such certificates. The found optimal solution can be used to confirm Gurobi's final primal bound. The log output contains the BB tree information, but it seems impossible to completely reconstruct Gurobi's BB tree or view the specific added cuts. Admittedly, using Gurobi means a certain trade-off between certifying the global optimality and efficiency of the computations. In our experiments, we verified the solutions of BB metods by comparing them with Lasserre's hierarchies where possible.

\subsection{Finding Feasible Solutions}
Getting feasible solutions as early as possible in the BB tree is essential for performance. Using an NLP solver on the problem is an obvious heuristic for finding feasible solutions, and is often performed internally by solvers. However, reformulating problems such that they can be solved by BB methods, e.g. by means of introducing auxiliary variables, can make result in a more complex program and hide some of the original structure. As a result, an NLP solver executed before any reformulation may have an easier time finding a feasible solutions even if the BB solver attempts the same.

Additionally, there may already exist domain-specific methods for finding heuristic solutions with significantly lower complexity than BB methods. For example, in robotics, there exist methods that work in real-time. These methods could be useful both for warm-starting and as runtime heuristics for BB.